\documentclass[11pt,a4paper]{article}
\usepackage[utf8]{inputenc}
\usepackage[T1]{fontenc}
\usepackage[spanish,es-tabla]{babel}
\usepackage{lmodern}
\usepackage{amsmath,amssymb,amsfonts}
\usepackage{geometry}
\geometry{top=2.1cm,bottom=2.2cm,left=2.2cm,right=2.2cm}
\usepackage{microtype}
\usepackage{booktabs,array,multirow,tabularx,colortbl}
\usepackage{enumitem}
\usepackage{graphicx}
\usepackage{xcolor}
\usepackage{tcolorbox}
\usepackage{tikz}
\usetikzlibrary{arrows.meta,positioning,shapes.geometric,calc,fit,backgrounds}
\usepackage{pgfplots}
\pgfplotsset{compat=1.18}
\usepackage{float}
\usepackage{caption}
\usepackage[hidelinks]{hyperref}
\usepackage{fancyhdr}
\usepackage{lastpage}

\definecolor{rejectred}{RGB}{190,55,55}
\definecolor{acceptblue}{RGB}{55,105,175}
\definecolor{altorange}{RGB}{220,125,45}
\definecolor{accentgreen}{RGB}{55,145,95}
\definecolor{softgray}{RGB}{245,246,248}
\definecolor{purplebox}{RGB}{115,80,155}

\pagestyle{fancy}
\fancyhf{}
\lhead{Estadística II -- Unidad 2}
\rhead{Pruebas de Hipótesis}
\cfoot{\thepage\ de \pageref{LastPage}}
\setlength{\headheight}{14pt}

\newtcolorbox{ideaclave}{
  colback=acceptblue!5,colframe=acceptblue!75!black,
  title=\textbf{Idea clave},fonttitle=\bfseries,arc=2mm,boxrule=.8pt}
\newtcolorbox{regla}{
  colback=accentgreen!6,colframe=accentgreen!75!black,
  title=\textbf{Regla para recordar},fonttitle=\bfseries,arc=2mm,boxrule=.8pt}
\newtcolorbox{errorfrecuente}{
  colback=rejectred!5,colframe=rejectred!75!black,
  title=\textbf{Error frecuente},fonttitle=\bfseries,arc=2mm,boxrule=.8pt}
\newtcolorbox{aclaracion}{
  colback=yellow!8,colframe=orange!80!black,
  title=\textbf{Aclaración de la guía},fonttitle=\bfseries,arc=2mm,boxrule=.8pt}
\newtcolorbox{oral}{
  colback=purplebox!5,colframe=purplebox!75!black,
  title=\textbf{Cómo explicarlo oralmente},fonttitle=\bfseries,arc=2mm,boxrule=.8pt}
\newtcolorbox{observargrafico}{
  colback=softgray,colframe=black!45,
  title=\textbf{Qué debo observar en este gráfico},fonttitle=\bfseries,arc=1mm,boxrule=.6pt}

\setlist[itemize]{leftmargin=1.4em,itemsep=3pt,topsep=3pt}
\setlist[enumerate]{leftmargin=1.7em,itemsep=4pt,topsep=4pt}
\setlength{\parindent}{0pt}
\setlength{\parskip}{5pt}
\captionsetup{font=small,labelfont=bf}

\title{\vspace{-1.2cm}\textbf{UNIDAD 2 -- PRUEBAS DE HIPÓTESIS}\\[4pt]
\large Apunte visual para comprender, resolver y exponer en el final oral}
\author{Estadística Aplicada II}
\date{}

\begin{document}
\maketitle
\tableofcontents
\newpage

\section{Mapa mental de la unidad}

Antes de memorizar fórmulas, conviene ver la lógica completa de una prueba de hipótesis:

\begin{center}
\begin{tikzpicture}[node distance=7mm and 5mm, every node/.style={font=\small},
 box/.style={draw,rounded corners,minimum height=8mm,align=center,fill=acceptblue!5},
 arr/.style={-{Latex[length=2.3mm]},thick}]
\node[box] (p1) {Afirmación sobre\\un parámetro};
\node[box,right=of p1] (p2) {Formular\\$H_0$ y $H_a$};
\node[box,right=of p2] (p3) {Elegir\\$\alpha$};
\node[box,below=of p3] (p4) {Elegir estadístico\\$Z$, $t$, etc.};
\node[box,left=of p4] (p5) {Determinar\\valor crítico};
\node[box,left=of p5] (p6) {Calcular\\estadístico};
\node[box,below=of p6] (p7) {Comparar con\\región de rechazo};
\node[box,right=of p7] (p8) {Rechazar o\\no rechazar $H_0$};
\node[box,right=of p8] (p9) {Interpretar en el\\contexto del problema};
\draw[arr] (p1)--(p2); \draw[arr] (p2)--(p3); \draw[arr] (p3)--(p4);
\draw[arr] (p4)--(p5); \draw[arr] (p5)--(p6); \draw[arr] (p6)--(p7);
\draw[arr] (p7)--(p8); \draw[arr] (p8)--(p9);
\end{tikzpicture}
\end{center}

\begin{ideaclave}
Una prueba de hipótesis no consiste solamente en calcular una fórmula. Consiste en construir una \textbf{regla de decisión} y después decidir si la evidencia muestral es suficientemente extrema como para rechazar la hipótesis nula.
\end{ideaclave}

\section{Fundamentos de las pruebas de hipótesis}

\subsection{Hipótesis y prueba de hipótesis}

\textbf{Hipótesis:} es una suposición sobre un parámetro de la población que requiere ser verificada mediante evidencia. Por ejemplo, afirmar que la comisión media mensual de una población de vendedores es $\mu=\$2000$.

\textbf{Prueba de hipótesis:} es un procedimiento estandarizado que permite decidir si una afirmación relativa a un parámetro es compatible o no con la evidencia obtenida de una muestra.

\begin{center}
\begin{tikzpicture}[node distance=8mm, box/.style={draw,rounded corners,align=center,minimum width=3.3cm,minimum height=9mm}, arr/.style={-{Latex},thick}]
\node[box,fill=acceptblue!5] (pop) {Población\\parámetro $\mu$};
\node[box,fill=softgray,right=of pop] (sample) {Muestra\\estadístico $\bar{x}$};
\node[box,fill=accentgreen!5,right=of sample] (decision) {Regla estadística\\de decisión};
\draw[arr] (pop)--node[above,font=\small]{seleccionamos}(sample);
\draw[arr] (sample)--node[above,font=\small]{calculamos}(decision);
\end{tikzpicture}
\end{center}

\subsection{Hipótesis nula e hipótesis alternativa}

\begin{itemize}
\item \textbf{Hipótesis Nula ($H_0$):} es el punto de partida. Se asume temporalmente para confrontarla con los datos. En la forma usada por la cátedra contiene la igualdad: $=$, $\le$ o $\ge$.
\item \textbf{Hipótesis Alternativa ($H_a$ o $H_1$):} expresa lo que se concluirá si existe evidencia suficiente para rechazar $H_0$. No lleva igualdad: $<$, $>$ o $\neq$.
\end{itemize}

\begin{regla}
El \textbf{no rechazo de $H_0$ no demuestra que $H_0$ sea verdadera}. Significa solamente que la muestra no aporta evidencia suficiente para rechazarla.
\end{regla}

\subsection{Hipótesis simple e hipótesis compuesta}

La guía distingue:
\begin{itemize}
\item \textbf{Hipótesis Simple:} especifica por completo un valor del parámetro, por ejemplo $\mu=20$.
\item \textbf{Hipótesis Compuesta:} comprende más de un valor posible del parámetro, por ejemplo $\mu\le20$ o $\mu>20$.
\end{itemize}

\subsection{Cómo elegir los signos}

\begin{table}[H]
\centering
\small
\begin{tabularx}{\textwidth}{>{\bfseries}p{3cm} X X}
\toprule
Situación & Planteamiento & Qué se busca \\
\midrule
Aumento / mejora & $H_0:\mu\le\mu_0$ vs. $H_a:\mu>\mu_0$ & Detectar valores significativamente mayores.\\
Disminución & $H_0:\mu\ge\mu_0$ vs. $H_a:\mu<\mu_0$ & Detectar valores significativamente menores.\\
Cambio / diferencia & $H_0:\mu=\mu_0$ vs. $H_a:\mu\neq\mu_0$ & Detectar una diferencia en cualquier dirección.\\
Garantía ``por lo menos'' & La garantía queda en $H_0$ y el incumplimiento en $H_a$ & Se supone válida la afirmación mientras la muestra no aporte evidencia en contra.\\
\bottomrule
\end{tabularx}
\end{table}

\begin{regla}
El signo de $H_a$ funciona como una flecha hacia la región de rechazo:
\[
H_a:\mu>\mu_0\Rightarrow\text{derecha},\qquad
H_a:\mu<\mu_0\Rightarrow\text{izquierda},\qquad
H_a:\mu\neq\mu_0\Rightarrow\text{dos colas}.
\]
\end{regla}

\subsection{Los tres gráficos que hay que saber de memoria}

\subsubsection{Prueba unilateral derecha}
\[
H_0:\mu\le\mu_0\qquad\text{vs.}\qquad H_a:\mu>\mu_0
\]
\begin{figure}[H]
\centering
\begin{tikzpicture}
\begin{axis}[width=11.2cm,height=5.2cm,domain=-3.8:3.8,samples=160,
 axis lines=middle,axis y line=none,ytick=\empty,xlabel={$z$},
 xtick={0,1.645},xticklabels={$0$,$z_{\alpha}=1.645$},ymin=0,ymax=.45,clip=false]
\addplot[draw=none,fill=acceptblue!18,domain=-3.8:1.645] {1/sqrt(2*pi)*exp(-x^2/2)} \closedcycle;
\addplot[draw=none,fill=rejectred!38,domain=1.645:3.8] {1/sqrt(2*pi)*exp(-x^2/2)} \closedcycle;
\addplot[very thick,acceptblue!80!black] {1/sqrt(2*pi)*exp(-x^2/2)};
\draw[dashed,thick] (axis cs:1.645,0)--(axis cs:1.645,.105);
\node at (axis cs:-.65,.17) {\small Región de no rechazo};
\node[align=center] at (axis cs:2.45,.075) {\small\bfseries Rechazo\\$\alpha=0.05$};
\end{axis}
\end{tikzpicture}
\caption{Unilateral derecha: se rechaza para valores grandes del estadístico.}
\end{figure}
\begin{observargrafico}
La región roja está donde apunta $H_a:\mu>\mu_0$. Para $\alpha=0.05$, el límite normal es aproximadamente $1.645$. Si $z_{calc}>1.645$, se rechaza $H_0$.
\end{observargrafico}

\subsubsection{Prueba unilateral izquierda}
\[
H_0:\mu\ge\mu_0\qquad\text{vs.}\qquad H_a:\mu<\mu_0
\]
\begin{figure}[H]
\centering
\begin{tikzpicture}
\begin{axis}[width=11.2cm,height=5.2cm,domain=-3.8:3.8,samples=160,
 axis lines=middle,axis y line=none,ytick=\empty,xlabel={$z$},
 xtick={-1.645,0},xticklabels={$-z_{\alpha}=-1.645$,$0$},ymin=0,ymax=.45,clip=false]
\addplot[draw=none,fill=rejectred!38,domain=-3.8:-1.645] {1/sqrt(2*pi)*exp(-x^2/2)} \closedcycle;
\addplot[draw=none,fill=acceptblue!18,domain=-1.645:3.8] {1/sqrt(2*pi)*exp(-x^2/2)} \closedcycle;
\addplot[very thick,acceptblue!80!black] {1/sqrt(2*pi)*exp(-x^2/2)};
\draw[dashed,thick] (axis cs:-1.645,0)--(axis cs:-1.645,.105);
\node[align=center] at (axis cs:-2.45,.075) {\small\bfseries Rechazo\\$\alpha=0.05$};
\node at (axis cs:.65,.17) {\small Región de no rechazo};
\end{axis}
\end{tikzpicture}
\caption{Unilateral izquierda: se rechaza para valores pequeños del estadístico.}
\end{figure}
\begin{observargrafico}
La cola de rechazo queda a la izquierda porque $H_a$ contiene $<$. En esa mitad de la escala normal los valores críticos son negativos.
\end{observargrafico}

\subsubsection{Prueba bilateral}
\[
H_0:\mu=\mu_0\qquad\text{vs.}\qquad H_a:\mu\neq\mu_0
\]
\begin{figure}[H]
\centering
\begin{tikzpicture}
\begin{axis}[width=11.8cm,height=5.4cm,domain=-3.8:3.8,samples=160,
 axis lines=middle,axis y line=none,ytick=\empty,xlabel={$z$},
 xtick={-1.96,0,1.96},xticklabels={$-1.96$,$0$,$1.96$},ymin=0,ymax=.45,clip=false]
\addplot[draw=none,fill=rejectred!38,domain=-3.8:-1.96] {1/sqrt(2*pi)*exp(-x^2/2)} \closedcycle;
\addplot[draw=none,fill=acceptblue!18,domain=-1.96:1.96] {1/sqrt(2*pi)*exp(-x^2/2)} \closedcycle;
\addplot[draw=none,fill=rejectred!38,domain=1.96:3.8] {1/sqrt(2*pi)*exp(-x^2/2)} \closedcycle;
\addplot[very thick,acceptblue!80!black] {1/sqrt(2*pi)*exp(-x^2/2)};
\draw[dashed,thick] (axis cs:-1.96,0)--(axis cs:-1.96,.06);
\draw[dashed,thick] (axis cs:1.96,0)--(axis cs:1.96,.06);
\node at (axis cs:0,.19) {\small\bfseries $1-\alpha=0.95$};
\node at (axis cs:-2.75,.055) {\small $\alpha/2=0.025$};
\node at (axis cs:2.75,.055) {\small $\alpha/2=0.025$};
\end{axis}
\end{tikzpicture}
\caption{Bilateral con $\alpha=0.05$: el riesgo se reparte entre ambas colas.}
\end{figure}
\begin{observargrafico}
En una prueba bilateral, $\alpha$ se divide en dos partes. Por eso, con $\alpha=0.05$, cada cola tiene $0.025$ y los valores críticos son $\pm1.96$, no $\pm1.65$.
\end{observargrafico}

\section{Nivel de significancia y errores de decisión}

\subsection{Nivel de significancia $\alpha$}

El \textbf{nivel de significancia} es la probabilidad fijada de cometer un \textbf{Error Tipo I}: rechazar $H_0$ siendo verdadera. La guía utiliza habitualmente $0.05$, $0.01$ y $0.10$.

\begin{table}[H]
\centering
\begin{tabular}{cl}
\toprule
$\alpha$ & Aplicación indicada en el apunte de estudio \\
\midrule
$0.05$ & Proyectos de investigación relacionados con consumidores.\\
$0.01$ & Control de calidad.\\
$0.10$ & Encuestas.\\
\bottomrule
\end{tabular}
\end{table}

\subsection{Error Tipo I, Error Tipo II y potencia}

\begin{table}[H]
\centering
\renewcommand{\arraystretch}{1.35}
\begin{tabular}{p{3.6cm}cc}
\toprule
 & \textbf{$H_0$ verdadera} & \textbf{$H_0$ falsa}\\
\midrule
\textbf{Rechazar $H_0$} & \cellcolor{rejectred!16}\textbf{Error Tipo I} $\alpha$ & \cellcolor{accentgreen!14}\textbf{Decisión correcta} $1-\beta$\\
\textbf{No rechazar $H_0$} & \cellcolor{acceptblue!12}\textbf{Decisión correcta} $1-\alpha$ & \cellcolor{altorange!16}\textbf{Error Tipo II} $\beta$\\
\bottomrule
\end{tabular}
\caption{Las cuatro posibilidades al decidir sobre $H_0$.}
\end{table}

\begin{itemize}
\item \textbf{Error Tipo I ($\alpha$):} rechazar $H_0$ cuando $H_0$ es verdadera.
\item \textbf{Error Tipo II ($\beta$):} no rechazar $H_0$ cuando $H_0$ es falsa.
\item \textbf{Potencia ($1-\beta$):} rechazar correctamente una $H_0$ falsa; es decir, detectar un cambio real.
\end{itemize}

\subsection{Ejemplo de la pintura: ver $\alpha$, $\beta$ y potencia juntos}

La guía parte de una pintura cuyo tiempo medio declarado es $20$ min. Se toman $n=36$ tableros, con $\sigma=2.4$ min, y se rechaza el producto si $\bar{x}>20.75$ min.

\[
\sigma_{\bar{x}}=\frac{2.4}{\sqrt{36}}=0.4\text{ min}
\]

\textbf{Si $H_0$ es verdadera, $\mu=20$:}
\[
z=\frac{20.75-20}{0.4}=1.875,
\qquad
\alpha=P(Z>1.875)\approx0.0304.
\]

\textbf{Si la media real es $\mu=21$:}
\[
z=\frac{20.75-21}{0.4}=-0.625,
\qquad
\beta=P(Z\le-0.625)\approx0.2660.
\]

Por lo tanto:
\[
1-\beta\approx0.7340.
\]

\begin{figure}[H]
\centering
\begin{tikzpicture}
\begin{axis}[width=13cm,height=6.3cm,xmin=18.3,xmax=22.6,ymin=0,ymax=1.08,
 axis lines=left,xlabel={Media muestral $\bar{x}$ (min)},ylabel={Densidad},
 xtick={20.75},xticklabels={$\bar{x}_c=20.75$},ytick=\empty,clip=false]
\addplot[thick,acceptblue] coordinates {(18.200,0.00004) (18.240,0.00006) (18.280,0.00010) (18.320,0.00015) (18.360,0.00022) (18.400,0.00033) (18.440,0.00050) (18.480,0.00073) (18.520,0.00106) (18.560,0.00153) (18.600,0.00218) (18.640,0.00308) (18.680,0.00431) (18.720,0.00596) (18.760,0.00817) (18.800,0.01108) (18.840,0.01488) (18.880,0.01979) (18.920,0.02605) (18.960,0.03396) (19.000,0.04382) (19.040,0.05599) (19.080,0.07082) (19.120,0.08869) (19.160,0.10996) (19.200,0.13498) (19.240,0.16404) (19.280,0.19738) (19.320,0.23512) (19.360,0.27730) (19.400,0.32379) (19.440,0.37432) (19.480,0.42842) (19.520,0.48547) (19.560,0.54463) (19.600,0.60493) (19.640,0.66521) (19.680,0.72423) (19.720,0.78063) (19.760,0.83306) (19.800,0.88016) (19.840,0.92068) (19.880,0.95347) (19.920,0.97761) (19.960,0.99238) (20.000,0.99736) (20.040,0.99238) (20.080,0.97761) (20.120,0.95347) (20.160,0.92068) (20.200,0.88016) (20.240,0.83306) (20.280,0.78063) (20.320,0.72423) (20.360,0.66521) (20.400,0.60493) (20.440,0.54463) (20.480,0.48547) (20.520,0.42842) (20.560,0.37432) (20.600,0.32379) (20.640,0.27730) (20.680,0.23512) (20.720,0.19738) (20.760,0.16404) (20.800,0.13498) (20.840,0.10996) (20.880,0.08869) (20.920,0.07082) (20.960,0.05599) (21.000,0.04382) (21.040,0.03396) (21.080,0.02605) (21.120,0.01979) (21.160,0.01488) (21.200,0.01108) (21.240,0.00817) (21.280,0.00596) (21.320,0.00431) (21.360,0.00308) (21.400,0.00218) (21.440,0.00153) (21.480,0.00106) (21.520,0.00073) (21.560,0.00050) (21.600,0.00033) (21.640,0.00022) (21.680,0.00015) (21.720,0.00010) (21.760,0.00006) (21.800,0.00004) (21.840,0.00003) (21.880,0.00002) (21.920,0.00001) (21.960,0.00001) (22.000,0.00000) (22.040,0.00000) (22.080,0.00000) (22.120,0.00000) (22.160,0.00000) (22.200,0.00000) (22.240,0.00000) (22.280,0.00000) (22.320,0.00000) (22.360,0.00000) (22.400,0.00000) (22.440,0.00000) (22.480,0.00000) (22.520,0.00000) (22.560,0.00000) (22.600,0.00000) (22.640,0.00000) (22.680,0.00000) (22.720,0.00000) (22.760,0.00000) (22.800,0.00000)};
\addplot[thick,altorange] coordinates {(18.200,0.00000) (18.240,0.00000) (18.280,0.00000) (18.320,0.00000) (18.360,0.00000) (18.400,0.00000) (18.440,0.00000) (18.480,0.00000) (18.520,0.00000) (18.560,0.00000) (18.600,0.00000) (18.640,0.00000) (18.680,0.00000) (18.720,0.00000) (18.760,0.00000) (18.800,0.00000) (18.840,0.00000) (18.880,0.00000) (18.920,0.00000) (18.960,0.00000) (19.000,0.00000) (19.040,0.00001) (19.080,0.00001) (19.120,0.00002) (19.160,0.00003) (19.200,0.00004) (19.240,0.00006) (19.280,0.00010) (19.320,0.00015) (19.360,0.00022) (19.400,0.00033) (19.440,0.00050) (19.480,0.00073) (19.520,0.00106) (19.560,0.00153) (19.600,0.00218) (19.640,0.00308) (19.680,0.00431) (19.720,0.00596) (19.760,0.00817) (19.800,0.01108) (19.840,0.01488) (19.880,0.01979) (19.920,0.02605) (19.960,0.03396) (20.000,0.04382) (20.040,0.05599) (20.080,0.07082) (20.120,0.08869) (20.160,0.10996) (20.200,0.13498) (20.240,0.16404) (20.280,0.19738) (20.320,0.23512) (20.360,0.27730) (20.400,0.32379) (20.440,0.37432) (20.480,0.42842) (20.520,0.48547) (20.560,0.54463) (20.600,0.60493) (20.640,0.66521) (20.680,0.72423) (20.720,0.78063) (20.760,0.83306) (20.800,0.88016) (20.840,0.92068) (20.880,0.95347) (20.920,0.97761) (20.960,0.99238) (21.000,0.99736) (21.040,0.99238) (21.080,0.97761) (21.120,0.95347) (21.160,0.92068) (21.200,0.88016) (21.240,0.83306) (21.280,0.78063) (21.320,0.72423) (21.360,0.66521) (21.400,0.60493) (21.440,0.54463) (21.480,0.48547) (21.520,0.42842) (21.560,0.37432) (21.600,0.32379) (21.640,0.27730) (21.680,0.23512) (21.720,0.19738) (21.760,0.16404) (21.800,0.13498) (21.840,0.10996) (21.880,0.08869) (21.920,0.07082) (21.960,0.05599) (22.000,0.04382) (22.040,0.03396) (22.080,0.02605) (22.120,0.01979) (22.160,0.01488) (22.200,0.01108) (22.240,0.00817) (22.280,0.00596) (22.320,0.00431) (22.360,0.00308) (22.400,0.00218) (22.440,0.00153) (22.480,0.00106) (22.520,0.00073) (22.560,0.00050) (22.600,0.00033) (22.640,0.00022) (22.680,0.00015) (22.720,0.00010) (22.760,0.00006) (22.800,0.00004)};
% alfa bajo H0
\addplot[draw=none,fill=rejectred!45,domain=20.75:22.6,samples=100] {1/(0.4*sqrt(2*pi))*exp(-((x-20)^2)/(2*0.4^2))} \closedcycle;
% beta bajo H1
\addplot[draw=none,fill=altorange!25,domain=18.3:20.75,samples=100] {1/(0.4*sqrt(2*pi))*exp(-((x-21)^2)/(2*0.4^2))} \closedcycle;
\draw[dashed,very thick] (axis cs:20.75,0)--(axis cs:20.75,1.0);
\node[acceptblue] at (axis cs:19.75,.93) {\small Bajo $H_0$};
\node[altorange] at (axis cs:21.28,.93) {\small Bajo $H_1$};
\node[rejectred!90!black,align=center] at (axis cs:20.98,.13) {\small $\alpha$\\Error I};
\node[altorange!90!black,align=center] at (axis cs:20.28,.50) {\small $\beta$\\Error II};
\node[accentgreen!70!black,align=center] at (axis cs:21.45,.44) {\small zona de potencia\\$1-\beta$};
\end{axis}
\end{tikzpicture}
\caption{Los dos errores se entienden mirando la misma frontera de decisión desde dos realidades distintas.}
\end{figure}

\begin{observargrafico}
La línea crítica $20.75$ no cambia. Si la realidad es $\mu=20$, el área que queda a la derecha es $\alpha$. Si la realidad es $\mu=21$, el área que queda a la izquierda es $\beta$ y el área a la derecha es la potencia $1-\beta$.
\end{observargrafico}

\begin{regla}
Para un criterio dado, reducir mucho $\alpha$ hace más difícil rechazar $H_0$ y puede aumentar $\beta$. La guía resume esta idea diciendo que no conviene elegir un $\alpha$ excesivamente pequeño sin considerar el Error Tipo II.
\end{regla}

\begin{figure}[H]
\centering
\begin{tikzpicture}
\begin{scope}[xshift=0cm]
\begin{axis}[width=6.4cm,height=4.3cm,xmin=-3,xmax=5,ymin=0,ymax=.43,axis lines=middle,axis y line=none,ytick=\empty,xtick={1.3},xticklabels={$c_1$},title={\small Frontera menos exigente},clip=false]
\addplot[thick,acceptblue,domain=-3:3,samples=120]{1/sqrt(2*pi)*exp(-x^2/2)};
\addplot[thick,altorange,domain=-1:5,samples=120]{1/sqrt(2*pi)*exp(-(x-2.2)^2/2)};
\addplot[draw=none,fill=rejectred!35,domain=1.3:4,samples=80]{1/sqrt(2*pi)*exp(-x^2/2)}\closedcycle;
\addplot[draw=none,fill=altorange!22,domain=-1:1.3,samples=80]{1/sqrt(2*pi)*exp(-(x-2.2)^2/2)}\closedcycle;
\draw[dashed](axis cs:1.3,0)--(axis cs:1.3,.29);
\node[rejectred!90!black] at(axis cs:1.7,.07){\scriptsize $\alpha$ mayor};
\node[altorange!90!black] at(axis cs:.5,.12){\scriptsize $\beta$ menor};
\end{axis}
\end{scope}
\begin{scope}[xshift=7.1cm]
\begin{axis}[width=6.4cm,height=4.3cm,xmin=-3,xmax=5,ymin=0,ymax=.43,axis lines=middle,axis y line=none,ytick=\empty,xtick={1.8},xticklabels={$c_2$},title={\small Frontera más exigente},clip=false]
\addplot[thick,acceptblue,domain=-3:3,samples=120]{1/sqrt(2*pi)*exp(-x^2/2)};
\addplot[thick,altorange,domain=-1:5,samples=120]{1/sqrt(2*pi)*exp(-(x-2.2)^2/2)};
\addplot[draw=none,fill=rejectred!35,domain=1.8:4,samples=80]{1/sqrt(2*pi)*exp(-x^2/2)}\closedcycle;
\addplot[draw=none,fill=altorange!22,domain=-1:1.8,samples=80]{1/sqrt(2*pi)*exp(-(x-2.2)^2/2)}\closedcycle;
\draw[dashed](axis cs:1.8,0)--(axis cs:1.8,.39);
\node[rejectred!90!black] at(axis cs:2.1,.055){\scriptsize $\alpha$ menor};
\node[altorange!90!black] at(axis cs:.8,.13){\scriptsize $\beta$ mayor};
\end{axis}
\end{scope}
\end{tikzpicture}
\caption{Con las distribuciones y el tamaño muestral fijos, desplazar la frontera para reducir $\alpha$ aumenta la región de no rechazo bajo la alternativa y, por tanto, $\beta$.}
\end{figure}

\section{Procedimiento general y elección del estadístico}

\subsection{Los cinco pasos de la cátedra}

\begin{enumerate}
\item Formular $H_0$ y una $H_a$ apropiada.
\item Especificar la probabilidad de Error Tipo I, es decir, el nivel $\alpha$; cuando corresponda, considerar también $\beta$.
\item Construir un criterio de decisión a partir de la distribución muestral de un estadístico apropiado.
\item Calcular con los datos el valor del estadístico.
\item Decidir si se rechaza o no se rechaza $H_0$ y expresar la conclusión en el contexto del problema.
\end{enumerate}

\begin{center}
\begin{tikzpicture}[node distance=6mm, every node/.style={font=\small},
 box/.style={draw,rounded corners,align=center,text width=6.1cm,minimum height=8mm,fill=softgray},
 yes/.style={draw,diamond,aspect=2,align=center,fill=yellow!10,text width=3.4cm},arr/.style={-{Latex},thick}]
\node[box] (a) {1. Leer qué afirmación se desea comprobar};
\node[box,below=of a] (b) {2. Formular $H_0$ y $H_a$};
\node[yes,below=of b] (c) {¿$H_a$ tiene $>$, $<$ o $\neq$?};
\node[box,below=of c] (d) {3. Fijar $\alpha$ y elegir $Z$, $t$ u otro estadístico};
\node[box,below=of d] (e) {4. Obtener valor crítico y calcular el estadístico};
\node[yes,below=of e] (f) {¿Cae en región de rechazo?};
\node[box,below left=8mm and -2mm of f,fill=rejectred!8] (g) {Sí: rechazar $H_0$};
\node[box,below right=8mm and -2mm of f,fill=acceptblue!8] (h) {No: no rechazar $H_0$};
\node[box,below=18mm of f] (i) {5. Interpretar en palabras del problema};
\draw[arr](a)--(b);\draw[arr](b)--(c);\draw[arr](c)--(d);\draw[arr](d)--(e);\draw[arr](e)--(f);
\draw[arr](f)--(g);\draw[arr](f)--(h);\draw[arr](g)--(i);\draw[arr](h)--(i);
\end{tikzpicture}
\end{center}

\subsection{Cuándo usar $Z$ y cuándo usar $t$}

Para el contraste de una media, la guía trabaja con el siguiente criterio práctico:

\begin{table}[H]
\centering
\small
\begin{tabularx}{\textwidth}{p{3.2cm}p{2.6cm}X}
\toprule
Situación & Estadístico & Idea \\
\midrule
$\sigma$ conocida & $Z$ & Se conoce la variabilidad poblacional.\\
Muestra grande ($n\ge30$), $\sigma$ desconocida & $Z$ usando $s$ & La guía usa la aproximación normal para muestras grandes.\\
Muestra pequeña, $\sigma$ desconocida y población normal & $t$ de Student & Se usa $s$ y $\nu=n-1$ grados de libertad.\\
\bottomrule
\end{tabularx}
\end{table}

\begin{center}
\begin{tikzpicture}[node distance=6mm and 3mm, every node/.style={font=\small},
 box/.style={draw,rounded corners,align=center,minimum height=8mm,text width=3.15cm},arr/.style={-{Latex},thick}]
\node[box,fill=softgray] (q1) {¿Se conoce $\sigma$?};
\node[box,fill=acceptblue!7,below left=of q1] (z1) {Sí $\Rightarrow$ usar $Z$};
\node[box,fill=softgray,below right=of q1] (q2) {No: ¿$n\ge30$?};
\node[box,fill=acceptblue!7,below left=of q2] (z2) {Sí $\Rightarrow$ usar $Z$ con $s$};
\node[box,fill=softgray,below right=of q2] (q3) {No: muestra pequeña y población normal};
\node[box,fill=altorange!8,below=of q3] (t) {Usar $t$ con $\nu=n-1$};
\draw[arr](q1)--node[left]{Sí}(z1);\draw[arr](q1)--node[right]{No}(q2);
\draw[arr](q2)--node[left]{Sí}(z2);\draw[arr](q2)--node[right]{No}(q3);\draw[arr](q3)--(t);
\end{tikzpicture}
\end{center}

\subsection{Por qué $t$ tiene colas más anchas}

\begin{figure}[H]
\centering
\begin{tikzpicture}
\begin{axis}[width=11.5cm,height=5.4cm,xmin=-4,xmax=4,ymin=0,ymax=.43,
 axis lines=left,xlabel={valor del estadístico},ylabel={densidad},ytick=\empty,legend style={draw=none,at={(.98,.98)},anchor=north east}]
\addplot[very thick,acceptblue,domain=-4:4,samples=160] {1/sqrt(2*pi)*exp(-x^2/2)};\addlegendentry{$Z$ normal}
\addplot[very thick,rejectred] coordinates {(-4.000,0.00512) (-3.920,0.00562) (-3.840,0.00616) (-3.760,0.00677) (-3.680,0.00744) (-3.600,0.00819) (-3.520,0.00902) (-3.440,0.00995) (-3.360,0.01098) (-3.280,0.01213) (-3.200,0.01341) (-3.120,0.01483) (-3.040,0.01643) (-2.960,0.01821) (-2.880,0.02019) (-2.800,0.02242) (-2.720,0.02490) (-2.640,0.02767) (-2.560,0.03077) (-2.480,0.03423) (-2.400,0.03809) (-2.320,0.04240) (-2.240,0.04720) (-2.160,0.05255) (-2.080,0.05849) (-2.000,0.06509) (-1.920,0.07240) (-1.840,0.08047) (-1.760,0.08937) (-1.680,0.09913) (-1.600,0.10982) (-1.520,0.12146) (-1.440,0.13407) (-1.360,0.14766) (-1.280,0.16220) (-1.200,0.17766) (-1.120,0.19395) (-1.040,0.21096) (-0.960,0.22852) (-0.880,0.24645) (-0.800,0.26449) (-0.720,0.28236) (-0.640,0.29974) (-0.560,0.31628) (-0.480,0.33162) (-0.400,0.34538) (-0.320,0.35721) (-0.240,0.36678) (-0.160,0.37384) (-0.080,0.37815) (0.000,0.37961) (0.080,0.37815) (0.160,0.37384) (0.240,0.36678) (0.320,0.35721) (0.400,0.34538) (0.480,0.33162) (0.560,0.31628) (0.640,0.29974) (0.720,0.28236) (0.800,0.26449) (0.880,0.24645) (0.960,0.22852) (1.040,0.21096) (1.120,0.19395) (1.200,0.17766) (1.280,0.16220) (1.360,0.14766) (1.440,0.13407) (1.520,0.12146) (1.600,0.10982) (1.680,0.09913) (1.760,0.08937) (1.840,0.08047) (1.920,0.07240) (2.000,0.06509) (2.080,0.05849) (2.160,0.05255) (2.240,0.04720) (2.320,0.04240) (2.400,0.03809) (2.480,0.03423) (2.560,0.03077) (2.640,0.02767) (2.720,0.02490) (2.800,0.02242) (2.880,0.02019) (2.960,0.01821) (3.040,0.01643) (3.120,0.01483) (3.200,0.01341) (3.280,0.01213) (3.360,0.01098) (3.440,0.00995) (3.520,0.00902) (3.600,0.00819) (3.680,0.00744) (3.760,0.00677) (3.840,0.00616) (3.920,0.00562) (4.000,0.00512)};\addlegendentry{$t$, 5 gl}
\addplot[thick,altorange,dashed] coordinates {(-4.000,0.00052) (-3.920,0.00065) (-3.840,0.00081) (-3.760,0.00100) (-3.680,0.00123) (-3.600,0.00151) (-3.520,0.00186) (-3.440,0.00229) (-3.360,0.00280) (-3.280,0.00342) (-3.200,0.00417) (-3.120,0.00508) (-3.040,0.00616) (-2.960,0.00745) (-2.880,0.00900) (-2.800,0.01083) (-2.720,0.01299) (-2.640,0.01553) (-2.560,0.01850) (-2.480,0.02197) (-2.400,0.02600) (-2.320,0.03065) (-2.240,0.03599) (-2.160,0.04210) (-2.080,0.04903) (-2.000,0.05685) (-1.920,0.06563) (-1.840,0.07542) (-1.760,0.08626) (-1.680,0.09818) (-1.600,0.11119) (-1.520,0.12528) (-1.440,0.14040) (-1.360,0.15652) (-1.280,0.17352) (-1.200,0.19129) (-1.120,0.20968) (-1.040,0.22849) (-0.960,0.24752) (-0.880,0.26652) (-0.800,0.28523) (-0.720,0.30336) (-0.640,0.32063) (-0.560,0.33674) (-0.480,0.35139) (-0.400,0.36432) (-0.320,0.37528) (-0.240,0.38404) (-0.160,0.39044) (-0.080,0.39433) (0.000,0.39563) (0.080,0.39433) (0.160,0.39044) (0.240,0.38404) (0.320,0.37528) (0.400,0.36432) (0.480,0.35139) (0.560,0.33674) (0.640,0.32063) (0.720,0.30336) (0.800,0.28523) (0.880,0.26652) (0.960,0.24752) (1.040,0.22849) (1.120,0.20968) (1.200,0.19129) (1.280,0.17352) (1.360,0.15652) (1.440,0.14040) (1.520,0.12528) (1.600,0.11119) (1.680,0.09818) (1.760,0.08626) (1.840,0.07542) (1.920,0.06563) (2.000,0.05685) (2.080,0.04903) (2.160,0.04210) (2.240,0.03599) (2.320,0.03065) (2.400,0.02600) (2.480,0.02197) (2.560,0.01850) (2.640,0.01553) (2.720,0.01299) (2.800,0.01083) (2.880,0.00900) (2.960,0.00745) (3.040,0.00616) (3.120,0.00508) (3.200,0.00417) (3.280,0.00342) (3.360,0.00280) (3.440,0.00229) (3.520,0.00186) (3.600,0.00151) (3.680,0.00123) (3.760,0.00100) (3.840,0.00081) (3.920,0.00065) (4.000,0.00052)};\addlegendentry{$t$, 30 gl}
\end{axis}
\end{tikzpicture}
\caption{Con pocos grados de libertad, $t$ es más ancha; al aumentar los grados de libertad se aproxima a $Z$.}
\end{figure}

\begin{ideaclave}
\[
\nu\uparrow \quad\Longrightarrow\quad t\text{ de Student}\longrightarrow Z.
\]
La distribución $t$ compensa la incertidumbre adicional que aparece al estimar $\sigma$ mediante $s$ con una muestra pequeña.
\end{ideaclave}

\subsection{Anatomía de las fórmulas}

\textbf{Una media, $Z$:}
\[
 z_{calc}=\frac{\overbrace{\bar{x}}^{\text{media muestral}}-\overbrace{\mu_0}^{\text{valor de }H_0}}
 {\underbrace{\sigma/\sqrt{n}}_{\text{error estándar}}}
\]

\textbf{Una media, $t$:}
\[
 t_{calc}=\frac{\overbrace{\bar{x}}^{\text{media muestral}}-\overbrace{\mu_0}^{\text{valor de }H_0}}
 {\underbrace{s/\sqrt{n}}_{\text{error estándar estimado}}},
 \qquad \nu=n-1.
\]

\begin{ideaclave}
El numerador mide \textbf{cuánto se aleja la muestra de lo que afirma $H_0$}. El denominador mide \textbf{cuánta variación muestral es esperable por azar}. Por eso el estadístico expresa el alejamiento en unidades de error estándar.
\end{ideaclave}

\subsection{Cómo buscar el valor crítico en la tabla normal de la cátedra}

La tabla usada en el material muestra el área comprendida \textbf{desde $0$ hasta $z$}. Por eso:

\textbf{Una cola con $\alpha=0.05$:}
\[
0.5-0.05=0.4500.
\]
Se busca aproximadamente $0.4500$ dentro de la tabla. El valor cercano $0.4505$ corresponde a fila $1.6$ y columna $0.05$:
\[
z_{\alpha}\approx1.65.
\]

\textbf{Dos colas con $\alpha=0.01$:}
\[
\frac{\alpha}{2}=0.005,
\qquad
0.5-0.005=0.4950.
\]
El área cercana a $0.4950$ lleva a:
\[
z_{\alpha/2}\approx2.58.
\]

\begin{regla}
Antes de buscar en la tabla, decidir si el contraste es de \textbf{una cola o dos colas}. En dos colas, primero se divide $\alpha$ por $2$.
\end{regla}

\subsection{Comparar en escala $z$ o en la escala real}

La misma regla de decisión puede expresarse de dos maneras. Por ejemplo, para una prueba bilateral:
\[
-z_{\alpha/2}\le z_{calc}\le z_{\alpha/2}
\]
es equivalente a exigir que la media muestral quede entre dos límites reales:
\[
L_I=\mu_0-z_{\alpha/2}\frac{\sigma}{\sqrt n},
\qquad
L_S=\mu_0+z_{\alpha/2}\frac{\sigma}{\sqrt n}.
\]
Si $\sigma$ se reemplaza por $s$ en el procedimiento de muestra grande usado por la guía, se emplea el error estándar estimado correspondiente.

\begin{ideaclave}
No son dos pruebas diferentes: una compara el \textbf{estadístico tipificado} con valores críticos y la otra compara $\bar{x}$ con los \textbf{límites críticos en las unidades originales}.
\end{ideaclave}

\subsection{Valores críticos que aparecen repetidamente en la unidad}

\begin{table}[H]
\centering
\begin{tabular}{lcc}
\toprule
Prueba normal & $\alpha$ & Valor crítico \\
\midrule
Unilateral derecha & $0.05$ & $+1.645\approx1.65$\\
Unilateral izquierda & $0.05$ & $-1.645\approx-1.65$\\
Bilateral & $0.05$ & $\pm1.96$\\
Unilateral & $0.01$ & $\pm2.326\approx\pm2.33$ según la dirección\\
Bilateral & $0.01$ & $\pm2.576\approx\pm2.58$\\
\bottomrule
\end{tabular}
\end{table}

\section{Hipótesis relativa a una media: ejemplos que hay que saber leer}

\subsection{Ejemplo de cola derecha: nuevo filamento}

Una fábrica tiene media histórica $\mu_0=10000$ h y $\sigma=500$ h. Con un nuevo filamento, una muestra de $n=40$ bombillas da $\bar{x}=10200$ h. Se desea comprobar si aumentó la duración con $\alpha=0.05$.

\[
H_0:\mu\le10000,
\qquad
H_a:\mu>10000.
\]
\[
z_{calc}=\frac{10200-10000}{500/\sqrt{40}}\approx2.53.
\]
Como $2.53>1.645$, se rechaza $H_0$.

\begin{figure}[H]
\centering
\begin{tikzpicture}
\begin{axis}[width=10.8cm,height=4.6cm,domain=-3.5:3.5,samples=140,axis lines=middle,axis y line=none,ytick=\empty,
 xtick={0,1.645,2.53},xticklabels={$0$,$1.645$,$2.53$},ymin=0,ymax=.44,clip=false]
\addplot[draw=none,fill=rejectred!30,domain=1.645:3.5]{1/sqrt(2*pi)*exp(-x^2/2)}\closedcycle;
\addplot[thick,acceptblue]{1/sqrt(2*pi)*exp(-x^2/2)};
\draw[dashed,thick] (axis cs:1.645,0)--(axis cs:1.645,.11);
\draw[very thick,altorange] (axis cs:2.53,0)--(axis cs:2.53,.06);
\node[altorange!80!black] at (axis cs:2.53,.095){\small $z_{calc}$};
\end{axis}
\end{tikzpicture}
\caption{El estadístico calculado cae en la cola derecha de rechazo.}
\end{figure}

\subsection{Ejemplo bilateral: tubos fluorescentes}

Datos de la guía: $n=100$, $\bar{x}=1570$ h, $s=120$ h, $\mu_0=1600$ h, $\alpha=0.05$.
\[
H_0:\mu=1600,
\qquad
H_a:\mu\neq1600.
\]
\[
z_{calc}=\frac{1570-1600}{120/\sqrt{100}}=-2.50.
\]
Con valores críticos $\pm1.96$, se rechaza $H_0$ porque $-2.50<-1.96$.

\begin{figure}[H]
\centering
\begin{tikzpicture}
\begin{axis}[width=11.3cm,height=4.7cm,domain=-3.5:3.5,samples=140,axis lines=middle,axis y line=none,ytick=\empty,
 xtick={-1.96,0,1.96},xticklabels={$-1.96$,$0$,$1.96$},ymin=0,ymax=.44,clip=false]
\addplot[draw=none,fill=rejectred!30,domain=-3.5:-1.96]{1/sqrt(2*pi)*exp(-x^2/2)}\closedcycle;
\addplot[draw=none,fill=rejectred!30,domain=1.96:3.5]{1/sqrt(2*pi)*exp(-x^2/2)}\closedcycle;
\addplot[thick,acceptblue]{1/sqrt(2*pi)*exp(-x^2/2)};
\draw[dashed] (axis cs:-1.96,0)--(axis cs:-1.96,.06);\draw[dashed](axis cs:1.96,0)--(axis cs:1.96,.06);
\draw[very thick,altorange] (axis cs:-2.5,0)--(axis cs:-2.5,.03);
\node[altorange!80!black,align=center] at(axis cs:-2.62,.085)
 {\small $z_{calc}$\\[-1pt]\scriptsize $-2.50$};
\end{axis}
\end{tikzpicture}
\caption{El valor $-2.50$ queda fuera de la zona central de no rechazo.}
\end{figure}

\subsection{Ejemplo unilateral izquierdo: neumáticos}

La afirmación es que la vida útil media es \emph{al menos} $28000$ km. Se obtiene $n=40$, $\bar{x}=27463$, $s=1348$ y $\alpha=0.01$.
\[
H_0:\mu\ge28000,
\qquad
H_a:\mu<28000.
\]
\[
z_{calc}=\frac{27463-28000}{1348/\sqrt{40}}\approx-2.52.
\]
El valor crítico es $-2.33$. Como $-2.52<-2.33$, se rechaza $H_0$.

\begin{figure}[H]
\centering
\begin{tikzpicture}
\begin{axis}[width=10.8cm,height=4.6cm,domain=-3.7:3.5,samples=140,axis lines=middle,axis y line=none,ytick=\empty,
 xtick={-2.33,0},xticklabels={$-2.33$,$0$},ymin=0,ymax=.44,clip=false]
\addplot[draw=none,fill=rejectred!30,domain=-3.7:-2.33]{1/sqrt(2*pi)*exp(-x^2/2)}\closedcycle;
\addplot[thick,acceptblue]{1/sqrt(2*pi)*exp(-x^2/2)};
\draw[dashed] (axis cs:-2.33,0)--(axis cs:-2.33,.027);
\draw[very thick,altorange] (axis cs:-2.52,0)--(axis cs:-2.52,.017);
\node[altorange!80!black,align=center] at(axis cs:-2.72,.085)
 {\small $z_{calc}$\\[-1pt]\scriptsize $-2.52$};
\end{axis}
\end{tikzpicture}
\caption{El estadístico muestral se encuentra dentro de la región crítica izquierda.}
\end{figure}

\subsection{Ejemplo $t$ con muestra pequeña: bombillas}

La guía trabaja con $n=8$, $\bar{x}=1070$ h, valor de referencia $1120$ h y una desviación de $125$ h utilizada en el cálculo. La alternativa es $\mu<1120$ y $\alpha=0.05$.
\[
t_{calc}=\frac{1070-1120}{125/\sqrt8}\approx-1.131,
\qquad \nu=7.
\]
El valor crítico es $t_{0.05,7}=-1.895$. Como $-1.131>-1.895$, \textbf{no se rechaza $H_0$}.

\begin{aclaracion}
En el enunciado de la guía, $125$ h aparece como desviación histórica y luego se utiliza dentro del estadístico $t$. Para seguir el procedimiento tal como se desarrolla en la materia, se conserva ese valor en el cálculo y se centra el estudio en la regla de decisión con $t$ y $7$ grados de libertad.
\end{aclaracion}

\begin{figure}[H]
\centering
\begin{tikzpicture}
\begin{axis}[width=10.8cm,height=4.7cm,xmin=-4,xmax=4,ymin=0,ymax=.40,axis lines=middle,axis y line=none,ytick=\empty,
 xtick={-1.895,0},xticklabels={$-1.895$,$0$},clip=false]
\addplot[draw=none,fill=rejectred!30,domain=-4:-1.895,samples=100]{0.384991*(1+x^2/7)^(-4)}\closedcycle;
\addplot[thick,acceptblue,domain=-4:4,samples=160]{0.384991*(1+x^2/7)^(-4)};
\draw[dashed] (axis cs:-1.895,0)--(axis cs:-1.895,.07);
\draw[very thick,altorange] (axis cs:-1.131,0)--(axis cs:-1.131,.20);
\node[altorange!80!black,align=center] at(axis cs:-1.02,.25)
 {\small $t_{calc}$\\[-1pt]\scriptsize $-1.131$};
\end{axis}
\end{tikzpicture}
\caption{El estadístico $-1.131$ queda a la derecha del límite $-1.895$: no entra en la región de rechazo.}
\end{figure}

\begin{oral}
En un ejercicio, conviene explicar primero \textbf{por qué la prueba es derecha, izquierda o bilateral}; después justificar \textbf{por qué se usa $Z$ o $t$}; y recién entonces calcular. El valor numérico final tiene sentido solamente cuando se compara con la región crítica correspondiente.
\end{oral}

\section{Curvas Características de Operación (C.O.)}

\subsection{Idea intuitiva}

La \textbf{Curva Característica de Operación} estudia el Error Tipo II para distintos valores posibles de la verdadera media. La guía define:
\[
L(\mu)=P(\text{no rechazar }H_0\mid \mu\text{ es la media real}).
\]

Cuando la media real pertenece a la alternativa, $L(\mu)$ representa $\beta$ para ese valor de $\mu$.

\begin{ideaclave}
Una prueba de hipótesis responde: \emph{¿qué decido con esta muestra?} La Curva C.O. responde otra pregunta: \emph{¿qué tan fácil o difícil sería que mi prueba detectara distintos valores reales de $\mu$?}
\end{ideaclave}

\subsection{Cola derecha: ejemplo de la pintura}

Con $\mu_0=20$, $\sigma=2.4$, $n=36$ y frontera $\bar{x}_c=20.75$, se no rechaza $H_0$ cuando $\bar{x}\le20.75$. Por tanto:
\[
L(\mu)=P(\bar{x}\le20.75\mid\mu)
      =\Phi\!\left(\frac{20.75-\mu}{0.4}\right).
\]

\begin{table}[H]
\centering
\begin{tabular}{ccc}
\toprule
Media real $\mu$ & $z=(20.75-\mu)/0.4$ & $L(\mu)$\\
\midrule
19.50 & 3.125 & 0.999 \\
20.00 & 1.875 & 0.970 \\
20.50 & 0.625 & 0.734 \\
20.75 & 0.000 & 0.500 \\
21.00 & -0.625 & 0.266 \\
21.50 & -1.875 & 0.030 \\
22.00 & -3.125 & 0.001 \\
\bottomrule
\end{tabular}
\caption{Valores seleccionados de la Curva C.O. de cola derecha del ejemplo de pintura.}
\end{table}

\begin{figure}[H]
\centering
\begin{tikzpicture}
\begin{axis}[width=11.5cm,height=5.5cm,xmin=19.4,xmax=22.05,ymin=0,ymax=1.03,
 axis lines=left,xlabel={Media real $\mu$},ylabel={$L(\mu)$},ytick={0,.5,1},xtick={19.5,20,20.75,21.5,22},grid=major,grid style={black!10}]
\addplot[very thick,rejectred] coordinates {(18.500,1.00000) (18.550,1.00000) (18.600,1.00000) (18.650,1.00000) (18.700,1.00000) (18.750,1.00000) (18.800,1.00000) (18.850,1.00000) (18.900,1.00000) (18.950,1.00000) (19.000,0.99999) (19.050,0.99999) (19.100,0.99998) (19.150,0.99997) (19.200,0.99995) (19.250,0.99991) (19.300,0.99986) (19.350,0.99977) (19.400,0.99963) (19.450,0.99942) (19.500,0.99911) (19.550,0.99865) (19.600,0.99798) (19.650,0.99702) (19.700,0.99567) (19.750,0.99379) (19.800,0.99123) (19.850,0.98778) (19.900,0.98321) (19.950,0.97725) (20.000,0.96960) (20.050,0.95994) (20.100,0.94792) (20.150,0.93319) (20.200,0.91543) (20.250,0.89435) (20.300,0.86971) (20.350,0.84134) (20.400,0.80921) (20.450,0.77337) (20.500,0.73401) (20.550,0.69146) (20.600,0.64617) (20.650,0.59871) (20.700,0.54974) (20.750,0.50000) (20.800,0.45026) (20.850,0.40129) (20.900,0.35383) (20.950,0.30854) (21.000,0.26599) (21.050,0.22663) (21.100,0.19079) (21.150,0.15866) (21.200,0.13029) (21.250,0.10565) (21.300,0.08457) (21.350,0.06681) (21.400,0.05208) (21.450,0.04006) (21.500,0.03040) (21.550,0.02275) (21.600,0.01679) (21.650,0.01222) (21.700,0.00877) (21.750,0.00621) (21.800,0.00433) (21.850,0.00298) (21.900,0.00202) (21.950,0.00135) (22.000,0.00089)};
\draw[dashed] (axis cs:20.75,0)--(axis cs:20.75,.5);
\draw[dashed] (axis cs:19.4,.5)--(axis cs:20.75,.5);
\node at(axis cs:21.45,.78){\small $n=36$};
\end{axis}
\end{tikzpicture}
\caption{Curva C.O. para una alternativa de cola derecha con frontera fija en $20.75$ min.}
\end{figure}
\begin{observargrafico}
A medida que la verdadera media aumenta, resulta cada vez menos probable que la media muestral permanezca por debajo de $20.75$. Por eso $L(\mu)$ disminuye: el Error Tipo II se vuelve menor para alternativas más alejadas.
\end{observargrafico}

\subsection{Cola izquierda: imagen espejo}

Para la alternativa contraria, la guía usa una frontera inferior $19.25$ min. Se no rechaza $H_0$ si $\bar{x}\ge19.25$:
\[
L(\mu)=P(\bar{x}\ge19.25\mid\mu).
\]

\begin{figure}[H]
\centering
\begin{tikzpicture}
\begin{axis}[width=11.5cm,height=5.5cm,xmin=18.0,xmax=22.0,ymin=0,ymax=1.03,
 axis lines=left,xlabel={Media real $\mu$},ylabel={$L(\mu)$},ytick={0,.5,1},xtick={18,19,20,21,22},legend style={draw=none,at={(.5,-.18)},anchor=north,legend columns=2}]
\addplot[very thick,rejectred] coordinates {(18.500,1.00000) (18.550,1.00000) (18.600,1.00000) (18.650,1.00000) (18.700,1.00000) (18.750,1.00000) (18.800,1.00000) (18.850,1.00000) (18.900,1.00000) (18.950,1.00000) (19.000,0.99999) (19.050,0.99999) (19.100,0.99998) (19.150,0.99997) (19.200,0.99995) (19.250,0.99991) (19.300,0.99986) (19.350,0.99977) (19.400,0.99963) (19.450,0.99942) (19.500,0.99911) (19.550,0.99865) (19.600,0.99798) (19.650,0.99702) (19.700,0.99567) (19.750,0.99379) (19.800,0.99123) (19.850,0.98778) (19.900,0.98321) (19.950,0.97725) (20.000,0.96960) (20.050,0.95994) (20.100,0.94792) (20.150,0.93319) (20.200,0.91543) (20.250,0.89435) (20.300,0.86971) (20.350,0.84134) (20.400,0.80921) (20.450,0.77337) (20.500,0.73401) (20.550,0.69146) (20.600,0.64617) (20.650,0.59871) (20.700,0.54974) (20.750,0.50000) (20.800,0.45026) (20.850,0.40129) (20.900,0.35383) (20.950,0.30854) (21.000,0.26599) (21.050,0.22663) (21.100,0.19079) (21.150,0.15866) (21.200,0.13029) (21.250,0.10565) (21.300,0.08457) (21.350,0.06681) (21.400,0.05208) (21.450,0.04006) (21.500,0.03040) (21.550,0.02275) (21.600,0.01679) (21.650,0.01222) (21.700,0.00877) (21.750,0.00621) (21.800,0.00433) (21.850,0.00298) (21.900,0.00202) (21.950,0.00135) (22.000,0.00089)};\addlegendentry{cola derecha}
\addplot[very thick,acceptblue,dashed] coordinates {(18.000,0.00089) (18.050,0.00135) (18.100,0.00202) (18.150,0.00298) (18.200,0.00433) (18.250,0.00621) (18.300,0.00877) (18.350,0.01222) (18.400,0.01679) (18.450,0.02275) (18.500,0.03040) (18.550,0.04006) (18.600,0.05208) (18.650,0.06681) (18.700,0.08457) (18.750,0.10565) (18.800,0.13029) (18.850,0.15866) (18.900,0.19079) (18.950,0.22663) (19.000,0.26599) (19.050,0.30854) (19.100,0.35383) (19.150,0.40129) (19.200,0.45026) (19.250,0.50000) (19.300,0.54974) (19.350,0.59871) (19.400,0.64617) (19.450,0.69146) (19.500,0.73401) (19.550,0.77337) (19.600,0.80921) (19.650,0.84134) (19.700,0.86971) (19.750,0.89435) (19.800,0.91543) (19.850,0.93319) (19.900,0.94792) (19.950,0.95994) (20.000,0.96960) (20.050,0.97725) (20.100,0.98321) (20.150,0.98778) (20.200,0.99123) (20.250,0.99379) (20.300,0.99567) (20.350,0.99702) (20.400,0.99798) (20.450,0.99865) (20.500,0.99911) (20.550,0.99942) (20.600,0.99963) (20.650,0.99977) (20.700,0.99986) (20.750,0.99991) (20.800,0.99995) (20.850,0.99997) (20.900,0.99998) (20.950,0.99999) (21.000,0.99999) (21.050,1.00000) (21.100,1.00000) (21.150,1.00000) (21.200,1.00000) (21.250,1.00000) (21.300,1.00000) (21.350,1.00000) (21.400,1.00000) (21.450,1.00000) (21.500,1.00000)};\addlegendentry{cola izquierda}
\end{axis}
\end{tikzpicture}
\caption{Las curvas de cola derecha e izquierda son imágenes especulares.}
\end{figure}

\subsection{Prueba bilateral}

Si la región de no rechazo queda entre $19.25$ y $20.75$:
\[
L(\mu)=P(19.25\le\bar{x}\le20.75\mid\mu).
\]

\begin{figure}[H]
\centering
\begin{tikzpicture}
\begin{axis}[width=11.5cm,height=5.5cm,xmin=18.5,xmax=21.5,ymin=0,ymax=1.02,
 axis lines=left,xlabel={Media real $\mu$},ylabel={$L(\mu)$},ytick={0,.5,1},xtick={18.5,19,19.5,20,20.5,21,21.5},grid=major,grid style={black!10}]
\addplot[very thick,rejectred] coordinates {(18.500,0.03040) (18.550,0.04006) (18.600,0.05208) (18.650,0.06681) (18.700,0.08457) (18.750,0.10565) (18.800,0.13029) (18.850,0.15865) (18.900,0.19079) (18.950,0.22662) (19.000,0.26598) (19.050,0.30853) (19.100,0.35381) (19.150,0.40126) (19.200,0.45021) (19.250,0.49991) (19.300,0.54959) (19.350,0.59847) (19.400,0.64580) (19.450,0.69089) (19.500,0.73313) (19.550,0.77202) (19.600,0.80719) (19.650,0.83836) (19.700,0.86537) (19.750,0.88814) (19.800,0.90666) (19.850,0.92097) (19.900,0.93113) (19.950,0.93719) (20.000,0.93921) (20.050,0.93719) (20.100,0.93113) (20.150,0.92097) (20.200,0.90666) (20.250,0.88814) (20.300,0.86537) (20.350,0.83836) (20.400,0.80719) (20.450,0.77202) (20.500,0.73313) (20.550,0.69089) (20.600,0.64580) (20.650,0.59847) (20.700,0.54959) (20.750,0.49991) (20.800,0.45021) (20.850,0.40126) (20.900,0.35381) (20.950,0.30853) (21.000,0.26598) (21.050,0.22662) (21.100,0.19079) (21.150,0.15865) (21.200,0.13029) (21.250,0.10565) (21.300,0.08457) (21.350,0.06681) (21.400,0.05208) (21.450,0.04006) (21.500,0.03040)};
\draw[dashed] (axis cs:20,0)--(axis cs:20,.939);
\node at(axis cs:20.18,.89){\small máximo cerca de $\mu_0$};
\end{axis}
\end{tikzpicture}
\caption{En dos colas, la probabilidad de no rechazo es máxima cerca de $\mu_0$ y disminuye hacia ambos lados.}
\end{figure}

\begin{observargrafico}
La forma bilateral es simétrica: si la verdadera media se aleja de $20$ tanto hacia abajo como hacia arriba, disminuye la probabilidad de permanecer dentro del intervalo de no rechazo.
\end{observargrafico}

\subsection{Efecto de aumentar el tamaño muestral en el análisis de la guía}

La fuente compara $n=36$ con $n=50$ manteniendo fija la frontera $\bar{x}_c=20.75$ min. Al aumentar $n$ disminuye el error estándar:
\[
\frac{\sigma}{\sqrt{n}}\downarrow.
\]

\begin{figure}[H]
\centering
\begin{tikzpicture}
\begin{axis}[width=11.7cm,height=5.5cm,xmin=19.5,xmax=22,ymin=0,ymax=1.03,
 axis lines=left,xlabel={Media real $\mu$},ylabel={$L(\mu)$},ytick={0,.5,1},xtick={19.5,20,20.5,21,21.5,22},legend style={draw=none,at={(.98,.96)},anchor=north east}]
\addplot[very thick,rejectred] coordinates {(18.500,1.00000) (18.550,1.00000) (18.600,1.00000) (18.650,1.00000) (18.700,1.00000) (18.750,1.00000) (18.800,1.00000) (18.850,1.00000) (18.900,1.00000) (18.950,1.00000) (19.000,0.99999) (19.050,0.99999) (19.100,0.99998) (19.150,0.99997) (19.200,0.99995) (19.250,0.99991) (19.300,0.99986) (19.350,0.99977) (19.400,0.99963) (19.450,0.99942) (19.500,0.99911) (19.550,0.99865) (19.600,0.99798) (19.650,0.99702) (19.700,0.99567) (19.750,0.99379) (19.800,0.99123) (19.850,0.98778) (19.900,0.98321) (19.950,0.97725) (20.000,0.96960) (20.050,0.95994) (20.100,0.94792) (20.150,0.93319) (20.200,0.91543) (20.250,0.89435) (20.300,0.86971) (20.350,0.84134) (20.400,0.80921) (20.450,0.77337) (20.500,0.73401) (20.550,0.69146) (20.600,0.64617) (20.650,0.59871) (20.700,0.54974) (20.750,0.50000) (20.800,0.45026) (20.850,0.40129) (20.900,0.35383) (20.950,0.30854) (21.000,0.26599) (21.050,0.22663) (21.100,0.19079) (21.150,0.15866) (21.200,0.13029) (21.250,0.10565) (21.300,0.08457) (21.350,0.06681) (21.400,0.05208) (21.450,0.04006) (21.500,0.03040) (21.550,0.02275) (21.600,0.01679) (21.650,0.01222) (21.700,0.00877) (21.750,0.00621) (21.800,0.00433) (21.850,0.00298) (21.900,0.00202) (21.950,0.00135) (22.000,0.00089)};\addlegendentry{$n=36$}
\addplot[very thick,acceptblue,dashed] coordinates {(19.500,0.99988) (19.550,0.99980) (19.600,0.99965) (19.650,0.99940) (19.700,0.99901) (19.750,0.99839) (19.800,0.99744) (19.850,0.99600) (19.900,0.99387) (19.950,0.99079) (20.000,0.98644) (20.050,0.98041) (20.100,0.97226) (20.150,0.96145) (20.200,0.94743) (20.250,0.92964) (20.300,0.90755) (20.350,0.88070) (20.400,0.84878) (20.450,0.81162) (20.500,0.76931) (20.550,0.72216) (20.600,0.67073) (20.650,0.61586) (20.700,0.55856) (20.750,0.50000) (20.800,0.44144) (20.850,0.38414) (20.900,0.32927) (20.950,0.27784) (21.000,0.23069) (21.050,0.18838) (21.100,0.15122) (21.150,0.11930) (21.200,0.09245) (21.250,0.07036) (21.300,0.05257) (21.350,0.03855) (21.400,0.02774) (21.450,0.01959) (21.500,0.01356) (21.550,0.00921) (21.600,0.00613) (21.650,0.00400) (21.700,0.00256) (21.750,0.00161) (21.800,0.00099) (21.850,0.00060) (21.900,0.00035) (21.950,0.00020) (22.000,0.00012)};\addlegendentry{$n=50$}
\end{axis}
\end{tikzpicture}
\caption{Con la misma frontera real, una muestra mayor produce una transición más abrupta de la Curva C.O.}
\end{figure}

\begin{regla}
En el análisis de la guía, una muestra mayor hace que la distribución de $\bar{x}$ sea más estrecha y que la prueba discrimine con mayor nitidez entre medias cercanas y medias claramente alejadas del criterio.
\end{regla}

\subsection{Forma general utilizada para las curvas de una cola}

La guía introduce la distancia estandarizada:
\[
d=\frac{|\mu-\mu_0|}{\sigma}.
\]
Para $\alpha=0.05$ de una cola y $z_\alpha\approx1.65$, presenta el Error Tipo II en función de $d$ y $n$ mediante la distribución normal. Escribiendo $\Phi$ para el área acumulada normal, el algoritmo de una cola puede condensarse como:
\[
\beta=L(d,n)=\Phi\!\left(z_\alpha-d\sqrt n\right).
\]
Para dos colas, con el valor crítico $z_{\alpha/2}$:
\[
\beta=\Phi\!\left(z_{\alpha/2}-d\sqrt n\right)
      -\Phi\!\left(-z_{\alpha/2}-d\sqrt n\right).
\]
El punto conceptual a recordar es:
\[
 d\sqrt{n}\ \text{crece}\quad\Rightarrow\quad \beta\ \text{disminuye}.
\]

Como ejemplo, para una prueba de una cola con $\mu_0=75.2$, $\mu=77$, $\sigma=3.6$, $n=15$:
\[
d=\frac{|77-75.2|}{3.6}=0.5,
\qquad
\beta\approx0.3873.
\]

\begin{oral}
La Curva C.O. es una forma de estudiar el Error Tipo II. Para cada valor real posible de $\mu$, indica la probabilidad de que la prueba no rechace $H_0$. Si el valor real se aleja en la dirección de la alternativa, esa probabilidad debería disminuir; por eso aumenta la capacidad de la prueba para detectar el cambio.
\end{oral}

\section{Hipótesis relativa a dos medias independientes}

\subsection{Qué cambia respecto de una sola media}

Ahora se comparan dos poblaciones independientes con medias $\mu_1$ y $\mu_2$. La hipótesis nula general de la guía es:
\[
H_0:\mu_1-\mu_2=\delta,
\]
donde $\delta$ es una constante, frecuentemente $0$.

Para muestras independientes:
\[
E(\bar{x}_1-\bar{x}_2)=\mu_1-\mu_2,
\]
\[
\sigma^2_{\bar{x}_1-\bar{x}_2}=\frac{\sigma_1^2}{n_1}+\frac{\sigma_2^2}{n_2}.
\]

\textbf{Con varianzas conocidas o muestras grandes:}
\[
 z_{calc}=\frac{(\bar{x}_1-\bar{x}_2)-\delta}
 {\sqrt{\dfrac{\sigma_1^2}{n_1}+\dfrac{\sigma_2^2}{n_2}}}.
\]
Para muestras grandes, la guía sustituye $\sigma_1,\sigma_2$ por $s_1,s_2$.

\begin{table}[H]
\centering
\begin{tabular}{cc}
\toprule
Hipótesis alternativa & Rechazar $H_0$ si\\
\midrule
$\mu_1-\mu_2<\delta$ & $z<-z_\alpha$\\
$\mu_1-\mu_2>\delta$ & $z>z_\alpha$\\
$\mu_1-\mu_2\neq\delta$ & $z<-z_{\alpha/2}$ o $z>z_{\alpha/2}$\\
\bottomrule
\end{tabular}
\end{table}

\subsection{Ejemplo $Z$: resistencia de conductores}

La guía compara $32$ alambres ordinarios ($\bar{x}_1=0.136$, $s_1=0.004$ ohms) con $32$ alambres de aleación ($\bar{x}_2=0.083$, $s_2=0.005$ ohms). Se quiere probar si la reducción supera $0.050$ ohms:
\[
H_0:\mu_1-\mu_2=0.050,
\qquad
H_a:\mu_1-\mu_2>0.050.
\]
\[
z=\frac{(0.136-0.083)-0.050}
{\sqrt{0.004^2/32+0.005^2/32}}\approx2.65.
\]
Como $2.65>1.65$, se rechaza $H_0$ y se apoya la afirmación de que la aleación reduce la resistencia en más de $0.050$ ohms.

\subsection{Muestras pequeñas: prueba $t$ bimuestral}

Si las muestras son pequeñas, se desconocen las varianzas poblacionales y se supone que las poblaciones son normales con varianzas iguales, la guía utiliza una $t$ bimuestral.

Varianza combinada:
\[
 s_p^2=\frac{(n_1-1)s_1^2+(n_2-1)s_2^2}{n_1+n_2-2}.
\]
Estadístico:
\[
 t_{calc}=\frac{(\bar{x}_1-\bar{x}_2)-\delta}
 {s_p\sqrt{\dfrac1{n_1}+\dfrac1{n_2}}},
 \qquad \nu=n_1+n_2-2.
\]

\subsection{Ejemplo $t$ bimuestral: fertilizante y trigo}

Datos de la guía:
\[
n_1=n_2=12,\quad
\bar{x}_1=0.280,\ s_1=0.022,\quad
\bar{x}_2=0.264,\ s_2=0.020.
\]
Se busca un incremento:
\[
H_0:\mu_1-\mu_2=0,
\qquad
H_a:\mu_1-\mu_2>0.
\]
Con $\nu=22$, el valor crítico es $t_{0.05}=1.717$. La varianza combinada es:
\[
s_p^2\approx0.000442,
\qquad s_p\approx0.02102.
\]
\[
t_{calc}=\frac{0.280-0.264}{0.02102\sqrt{1/12+1/12}}\approx1.849.
\]
Como $1.849>1.717$, se rechaza $H_0$ y se concluye que el fertilizante produce un incremento significativo.

\begin{figure}[H]
\centering
\begin{tikzpicture}
\begin{axis}[width=10.5cm,height=4.5cm,xmin=-4,xmax=4,ymin=0,ymax=.40,axis lines=middle,axis y line=none,ytick=\empty,
 xtick={0,1.717},xticklabels={$0$,$1.717$},clip=false]
\addplot[draw=none,fill=rejectred!28,domain=1.717:4,samples=100]{0.394*(1+x^2/22)^(-11.5)}\closedcycle;
\addplot[thick,acceptblue,domain=-4:4,samples=160]{0.394*(1+x^2/22)^(-11.5)};
\draw[dashed](axis cs:1.717,0)--(axis cs:1.717,.09);
\draw[very thick,altorange](axis cs:1.849,0)--(axis cs:1.849,.075);
\node[altorange!80!black,align=center] at(axis cs:2.08,.105)
 {\small $t_{calc}$\\[-1pt]\scriptsize $1.849$};
\end{axis}
\end{tikzpicture}
\caption{El valor $t=1.849$ supera el crítico $1.717$ y entra en la cola de rechazo.}
\end{figure}

\subsection{Error Tipo II para dos medias}

La guía adapta las Curvas C.O. utilizando la diferencia alternativa real $\delta'=\mu_1-\mu_2$ y define:
\[
d=\frac{|\delta-\delta'|}{\sqrt{\sigma_1^2+\sigma_2^2}}.
\]
Si $n_1\neq n_2$, el tamaño muestral equivalente que aparece en la guía es:
\[
 n=\frac{(\sigma_1^2+\sigma_2^2)}
 {\dfrac{\sigma_1^2}{n_1}+\dfrac{\sigma_2^2}{n_2}}.
\]
En el ejemplo de estaturas de la fuente, para $\delta'=0.02$ m se obtiene $\beta=0.5549$.

\section{Muestras apareadas}

\subsection{Independientes vs. apareadas}

\begin{center}
\begin{tikzpicture}[every node/.style={font=\small},node distance=6mm and 14mm,
 box/.style={draw,rounded corners,align=center,minimum width=4cm,minimum height=9mm}]
\node[box,fill=acceptblue!6] (i1) {Muestra independiente 1\\$x_{11},x_{12},\ldots$};
\node[box,fill=acceptblue!6,right=of i1] (i2) {Muestra independiente 2\\$x_{21},x_{22},\ldots$};
\node[below=4mm of $(i1)!0.5!(i2)$] {\textbf{No existe correspondencia uno a uno}};
\node[box,fill=altorange!7,below=17mm of i1] (p1) {Unidad 1\\Antes $\leftrightarrow$ Después};
\node[box,fill=altorange!7,right=of p1] (p2) {Unidad 2\\Antes $\leftrightarrow$ Después};
\node[below=4mm of $(p1)!0.5!(p2)$] {\textbf{Cada observación tiene su pareja}};
\end{tikzpicture}
\end{center}

Cuando las observaciones vienen en pares, no se analizan como dos muestras independientes. Se define para cada par:
\[
d_i=x_{\text{Antes},i}-x_{\text{Después},i}.
\]
El problema se transforma en una prueba sobre \textbf{una sola media de diferencias}:
\[
\bar d=\frac{\sum d_i}{n},
\qquad
s_d=\sqrt{\frac{\sum(d_i-\bar d)^2}{n-1}},
\]
\[
 t_{calc}=\frac{\bar d-\mu_{d,0}}{s_d/\sqrt n},
 \qquad \nu=n-1.
\]

\begin{center}
\begin{tikzpicture}[node distance=9mm, every node/.style={font=\small},
 box/.style={draw,rounded corners,align=center,minimum width=4.2cm,minimum height=9mm},arr/.style={-{Latex},thick}]
\node[box,fill=altorange!7] (a) {Dos mediciones relacionadas};
\node[box,fill=softgray,right=of a] (b) {Calcular $d_i$};
\node[box,fill=accentgreen!7,right=of b] (c) {Prueba sobre $\mu_d$};
\draw[arr](a)--(b);\draw[arr](b)--(c);
\end{tikzpicture}
\end{center}

\subsection{Ejemplo: programa de seguridad}

La guía registra horas-hombre perdidas antes y después en diez plantas:

\begin{table}[H]
\centering
\small
\begin{tabular}{c*{10}{c}}
\toprule
Planta &1&2&3&4&5&6&7&8&9&10\\
\midrule
Antes &45&73&46&124&33&57&83&34&26&17\\
Después&36&60&44&119&35&51&77&29&24&11\\
$d_i$ &9&13&2&5&-2&6&6&5&2&6\\
\bottomrule
\end{tabular}
\end{table}

Se desea verificar si el programa reduce las pérdidas. Como $d_i=\text{Antes}-\text{Después}$, una reducción eficaz produce diferencias positivas:
\[
H_0:\mu_d=0,
\qquad
H_a:\mu_d>0.
\]
\[
\bar d=5.2,
\qquad s_d\approx4.077,
\qquad
 t_{calc}=\frac{5.2}{4.077/\sqrt{10}}\approx4.033.
\]
Con $\nu=9$ y $\alpha=0.05$, $t_{crit}=1.833$. Como $4.033>1.833$, se rechaza $H_0$: el sistema de seguridad es eficaz.

\begin{errorfrecuente}
No aplicar una prueba de muestras independientes a datos del tipo ``antes y después''. En datos apareados, la información esencial está en las diferencias de cada pareja.
\end{errorfrecuente}

\section{Síntesis final para el examen oral}

\begin{table}[H]
\centering
\small
\begin{tabularx}{\textwidth}{p{3.4cm}X}
\toprule
Pregunta que puede aparecer & Respuesta que debe estar clara \\
\midrule
¿Qué es una prueba de hipótesis? & Una regla de inferencia para contrastar una afirmación sobre un parámetro con evidencia muestral.\\
¿Qué es $H_0$? & La hipótesis de partida que se somete a prueba. No rechazarla no demuestra que sea verdadera.\\
¿Qué indica $H_a$? & Qué conclusión se busca y, mediante su signo, hacia qué cola se dirige la región crítica.\\
¿Qué es $\alpha$? & Probabilidad de Error Tipo I: rechazar $H_0$ siendo verdadera.\\
¿Qué es $\beta$? & Probabilidad de Error Tipo II: no rechazar $H_0$ siendo falsa para una alternativa concreta.\\
¿Qué es potencia? & $1-\beta$: probabilidad de detectar correctamente una alternativa real.\\
¿Cuándo uso $t$? & En el criterio de la guía: muestra pequeña, $\sigma$ desconocida y población normal; se usan $\nu=n-1$ gl.\\
¿Qué es una Curva C.O.? & La función $L(\mu)$ que muestra la probabilidad de no rechazo de $H_0$ para distintos valores reales del parámetro.\\
¿Independientes o apareadas? & Independientes: grupos sin emparejamiento. Apareadas: cada unidad tiene dos mediciones y se analiza $d_i$.\\
\bottomrule
\end{tabularx}
\end{table}

\begin{oral}
Una exposición segura sigue siempre el mismo hilo: \textbf{defino el concepto, planteo $H_0$ y $H_a$, justifico una o dos colas, explico el nivel $\alpha$, elijo el estadístico, ubico el valor crítico, comparo el estadístico calculado y termino interpretando la decisión en el contexto del problema}. Si me preguntan por la calidad de la prueba, conecto $\beta$, potencia y Curvas C.O.
\end{oral}

\section*{Fuentes de cátedra utilizadas}
\addcontentsline{toc}{section}{Fuentes de cátedra utilizadas}
\begin{itemize}
\item \textit{Capitulo 2 - Pruebas de hipotesis} (material de estudio del usuario).
\item \textit{2-hipotesis.docx} -- Cátedra Estadística Aplicada II, Universidad de Mendoza.
\end{itemize}

\end{document}
